\cotefile{157-4}
\scanpage{243}
\mspage{117}
% archive sheet number: 242

\uline{Loc 4} (Axiome du carré cocartésien) Soit $X$ dans \unsure{$X$|$\mathcal{X}$}, $X_1$, $X_2$ deux ouverts de réunion $X$, $X_0=X_1\cap X_2$, d'où un diagramme cocartésien dans $\mathrm{Cat}$

\diagram{\begin{tikzcd}[ampersand replacement=\&,column sep=small,row sep=small]
 X_0 \arrow[r,"i_1"] \arrow[d,"i_2"'] \& X_1 \arrow[d,"j_2"] \\
 X_2 \arrow[r,"j_1"] \& \unsure{x|X}
\end{tikzcd}}

On a alors $i_1\in W\iff j_1\in W$ (et symétriquement $i_2\in W\iff j_2\in W$)

% slanted
\margin{Il y aura un Loc (5) \deleted{\unsure{plus}} ``basique'' plus bas, p.~124}

\uline{Loc} \textcircled{5} (Axiomes des $W$-fibrations) : \deleted{\unsure{la variante}}\inserted{il y a des} \deleted{\unsure{plus f}} variantes \textcircled{a} a$'$) \textcircled{b} b$'$) b$''$) c, c$'$, cf \uline{\unsure{XII}} p.~53-59. La forme la plus forte \textcircled{a} est celle-ci :

Soit $f:X\longrightarrow Y$ dans $\mathrm{Cat}$ \deleted{telle que pour toute}\inserted{une $W$-fibration} \deleted{\unsure{flèche} $Y''\to Y'$ dans} \unsure{au sens ancien} i.e. telle que pour toute $Y''\to Y'$ dans $\mathrm{Cat}_{/Y}$, avec $Y',Y''$ contractiles, \deleted{$X''=X\times_Y Y''$} la flèche $X''\to X'$ déduite par le chgt de base $X\to Y$ soit dans $W$. Alors $f$ est dans $\mathrm{Fib}_W$, i.e. le foncteur de chgt.\ de base \unsure{sur $f$} $\mathrm{Cat}_{/Y}\to\mathrm{Cat}_{/X}$ transforme flèches dans $W$ en flèches dans $W$.

\deleted{\uline{Loc \unsure{6|8}} (axiome de connexité) : \unsure{pour que} $\pi_0(f)$ \unsure{bijectif}}

% slanted
\margin{\deleted{Il y \unsure{avait} un Loc 6 concernant \unsure{les|des} \illegible{1} limites inductives}}

\uline{Loc} \deleted{\unsure{6|5}}7 (Axiome des limites inductives) $W$ est stable par limites inductives filtrantes

\inserted{forme renforcée de Loc 6}

\uline{Loc \deleted{\unsure{6|8}}7 bis}) (L'Accessibilité) $W$ est \deleted{accessible}\inserted{stable par $\mathrm{Lim}$} filtrantes \uline{et} accessible, i.e. il existe

\margin{\unsure{\deleted{la}|la|le}}

\uline{NB} \deleted{la conjonction de Loc 6 et Loc 7 \unsure{équivaut} à la \unsure{forme}} une petite partie \deleted{\unsure{des}} $W_{\unsure{\pi|0}}$ de $\mathrm{Fl}(\unsure{\mathrm{Cat}|\mathit{Cat}})$, telle que $W$ soit la plus petite partie de $\mathrm{Fl}(\mathrm{Cat})$ stable par \uline{lim} filtrantes, contenant $W_\pi$.
\scanpage{244}
\mspage{118}
% archive sheet number: 243

% slanted
\margin{J'ai changé la numérotation, pour tenir compte d'un Loc (7) plus bas, qui est encore un axiome de stabilité.}

\uline{Loc \deleted{\unsure{7}}8} (Axiome de connexité) $f\in W\Rightarrow\pi_0(f)$ bijectif

\uline{NB} Tous ces axiomes, sauf le dernier (et un \unsure{chouïa|chouia} de l'avant-dernier) \unsure{sont} des axiomes de stabilité --- ils disent tous que $W$ est ``grand'' : \deleted{\unsure{S}}il contient certaines flèches, et s'il contient telles flèches, il contient telles autres.

Loc \deleted{\unsure{7}}8, et dans une certaine mesure Loc \deleted{\unsure{6}}7 bis, sont des axiomes ``limitatifs'' sur la taille de $W$ : Loc 7 nous dit que $W\subset W_0$ ($=$ ens des $f\in\mathrm{Fl}\,\mathrm{Cat}$ telles que $\pi_0(f)$ bij.), Loc \deleted{\unsure{6}}7 bis nous dit l'existence d'un cardinal tel que $W$ soit égal à la \inserted{$\unsure{L|k}$-}\unsure{enveloppe} de $W^\pi$ (donc $W\subset\unsure{\mathrm{l}\text{-}\mathrm{env}|\mathrm{L}\text{-}\mathrm{env}}(W^\pi)$).

J'ai \unsure{essayé|envisagé} de ranger les axiomes par ordre d'importance pratique. J'ai fini par \unsure{rejeter|regarder} l'axiome de connexité à la fin, car dans les développements théoriques sur les localiseurs généraux, il n'intervient pratiquement jamais.

Mais je rappelle (\uline{XII} p.~\unsure{10|20}, Th.~2) que si $W$ satisfait $L(1,2,3,\unsure{6|5})$, et si $L(\deleted{\unsure{7}}8)$ n'est \uline{pas} satisfait, \deleted{\unsure{il suffit à 5 b}} alors $W=\mathrm{Fl}(\mathrm{Cat})$ ; \deleted{\unsure{donc}} en présence de $L(1,2,3,5)$, $L(\deleted{\unsure{7}}8)$ \uline{équivaut} à $W\ne\mathrm{Fl}(\mathrm{Cat})$ tout entier.
\scanpage{245}
\mspage{119}
% archive sheet number: 244

Les axiomes qui ont servi pour 8.4. $1^\circ,3^\circ$ \deleted{\unsure{et 2°}}\inserted{et 8.5.b)} sont
\[
 L(1\,\mathrm{bis},2,3)\unsure{'.|/.|!}
\]

Ceux qui sont nécessaires pour 8.5. \unsure{a)} \deleted{\unsure{et}} \deleted{8.5} sont \deleted{$\unsure{L(1)}$}
\[
 L(1\,\mathrm{bis},2,3,4).
\]

Enfin, ceux nécessaires pour 8.4. $2^\circ$ sont
\[
 L(1\,\mathrm{bis},2,3,4,\deleted{\unsure{6}}7)
\]
et \unsure{ils} impliquent la \uline{saturation forte de $W$} --- (cf \S{} suivant).

$\backslash$

Il me faut réparer un oubli, ayant traité \deleted{par} par par le mépris (!) les questions de composantes connexes, jusqu'à présent.

\uline{Proposition 9.1}. Supposons que $W$ satisfait Loc $(1,\inserted{2,3,}4)$ sans plus. \deleted{(axiome de saturation, axiome du carré cocartésien).} Soient $f_i:X_i\unsure{-|\to}Y_i$ $(i\in I)$ une famille \uline{finie} de flèches de $\mathrm{Cat}$. Pour que $f=\coprod f_i:\coprod X_i\to\coprod Y_i$ soit \unsure{dans} $W$, il \deleted{\unsure{faut}} s. que $\forall i\in I$, on ait $f_i\in W$. Si $W$ satisfait de plus : Loc 6 (axiome des limites inductives), alors cette conclusion reste valable sans condition de finitude sur $I$.

\uline{Dém}. Le cas $I$ fini se traite par réc. \unsure{$n$|$u$}$=\mathrm{card}\,I$, on est ramené aussitôt au cas où $n=2$. * On voit d'abord \deleted{\unsure{sur}}\inserted{de}

% slanted
\margin{Le cas infini \unsure{s'y ramène via|s'y ramènera,} l'hyp. Loc (6), \unsure{qui en est la limite|puisque c'est une limite}.}
\scanpage{246}
\mspage{120}
% archive sheet number: 245

\uline{Corollaire 9.2}. \deleted{\unsure{Supposons|Réciproquement}} Sous l'hyp. Loc $(1,4)$, soit $f:X\unsure{\to|-}Y$, \deleted{\unsure{soit|et}}\inserted{\unsure{$\in$|une}} flèche de $\mathrm{Cat}$, et $Z$ objet de $\mathrm{Cat}$. Alors $f\in W$ ssi $f\amalg\mathrm{id}_Z:X\amalg Z-Y\amalg Z$ est $\in W$.

Immédiat sur le carré cocartésien, où \deleted{\unsure{est imm. ouv.}} $\mathrm{inj}_0^X$ est imm. ouv. de Dwyer

\diagram{\begin{tikzcd}[ampersand replacement=\&,column sep=large]
X \arrow[r,"f",no head] \arrow[d,"\mathrm{inj}_0^X"'] \& Y \arrow[d,"\mathrm{inj}_0^Y"] \\
X\amalg Z \arrow[r,"{f'=f\amalg\mathrm{id}_Z}"',no head] \& Y\amalg Z
\end{tikzcd}}

qui montre que $f\in W\unsure{\iff|\Longrightarrow} f'\in W$.

\unsure{:} En itérant ce résultat, on trouve, si
\[
 f_1:X_1-Y_1,\quad f_2:X_2-Y_2
\]
sont donnés, alors
\[
 f_1\in W,\ f_2\in W\Longrightarrow f_1\amalg f_2\in W.
\]
en regardant le diagramme \deleted{\unsure{bi}}\inserted{de deux carrés} cocartésiens : \deleted{\unsure{d'immersions ouvertes}}

\diagram{\begin{tikzcd}[ampersand replacement=\&,column sep=huge,row sep=large]
X_1 \arrow[r,"f_1"] \arrow[d] \& Y_1 \arrow[d] \\
X_1\amalg X_2 \arrow[r,"{\underbrace{f'_1=f_1\amalg\mathrm{id}_{X_2}}}"] \arrow[d,"{f'_2=\mathrm{id}_{X_1}\amalg f_2}"] \& Y_1\amalg X_2 \arrow[d,"{g_2=\mathrm{id}_{Y_1}\amalg f_2}"] \\
X_1\amalg Y_2 \arrow[r,"{g_1=\deleted{\unsure{f}}f_1\amalg\mathrm{id}_{Y_2}}"'] \& Y_1\amalg Y_2
\end{tikzcd}}

Si $f_1$ et $f_2$ sont dans $W$, il en est de $=$ de $f'_1,f'_2$, \deleted{\unsure{par 9.2}}\inserted{$g_1,g_2$}, \deleted{et si on prouve que}\inserted{donc aussi de $f_1\amalg f_2=g_2f'_1$.}

\deleted{chaque carré est $W$-cocartésien, on en conclut \inserted{\unsure{(cf juste hyp. sur $f_1,f_2$,|sans hyp. sur $f_1,f_2$,} par Loc $(4)$)} que $g_1,g_2\in W$, donc aussi $f_1\amalg f_2=g_2f'_1\in W$. (Mais je ne vois \unsure{aucun|aucune}}
\scanpage{247}
\mspage{121}
% archive sheet number: 246

\deleted{moyen de prouver \illegible{1} que ce carré est $W$-cocartésien --- donc finalement je ne suis pas capable de prouver que $f_1,f_2\in W\Rightarrow f_1\amalg f_2\in W$, sous \uline{les seules hypothèses Loc $(1,2,3,4)$ !}}

Mais je ne vois aucun moyen de prouver la réciproque, que $f_1\amalg f_2\in W\Rightarrow f_1,f_2\in W$. Grâce à \deleted{\unsure{9.1}} ce qui a été prouvé, on est ramené au cas où $f_1,f_2$ sont des imm. ouvertes de Dwyer, mais il ne semble pas que cela nous avance d'un cheveu ! \deleted{Ce qui} En somme, il faut savoir que les applications canoniques $Y_i\hookrightarrow\coprod Y_i$ sont des $\mathrm{Fib}_W$.

\deleted{\unsure{Comme}} \unsure{comme} $Y_i$ est un \unsure{summand|sommand} direct, on est ramené au cas d'une immersion can. $Y_1\hookrightarrow Y_1\amalg Y_2=Y$ : si $f:X\unsure{-|\to}Y$ est dans $W$, en est-il de $=$ de $X_1=X\unsure{\mathbin{\alpha}_Y|\times_Y}Y_1\Rightarrow\deleted{\unsure{X}}Y_1$ ?

C'est une propriété qui, bien sûr, \deleted{\unsure{serait}}\inserted{est} conséquence de l'ancien $W(4a)$\inserted{(\uline{XII} p.~53) ou seulement} $W(4a')$, donc de l'actuel Loc $(5a')$. Il suffit $=$ de la forme $W(4c')=\mathrm{Loc}\,(4c')$ la plus faible de toutes, comme on voit en décomposant $X_{\unsure{1}}-Y$ en $X\xrightarrow{i}X'\xrightarrow{p}Y$, avec $i\in W$ \deleted{\unsure{cof}} (plus précisément, $i'_1:X_1-X'_{\unsure{2|1}}$ et $i'_2:X_2-X'_2$ \deleted{\unsure{et}} dans $W$ !), et \deleted{$p\in \unsure{P}$} $p\in\mathrm{Parf}_W$. Par Loc $(4c')$\inserted{($f\in W$ \unsure{donc $p$ par saturation|aussi $p$ par saturation})}, les fibres de $p$ sont $W$-asphériques, ce qui implique que $p$ est alors $W$-\unsure{univ|universel}, \unsure{donc|avec} $p_1:X'_1-Y_1$ est dans $W$.
\scanpage{248}
\mspage{122}
% archive sheet number: 247

\unsure{Donc|Bon} notons pour acquit de conscience

\uline{Corollaire \unsure{9.3|5.3}}. Si $W$ satisfait $\mathrm{Loc}\,(1,2,3,\deleted{\unsure{4,}}5c')$ (où Loc $(\unsure{6|5}c')$ est la variante la plus faible de toutes de l'axiome des fibrations), alors \deleted{\unsure{pour t}} pour tte famille $(f_i:X_i\unsure{-|\to}Y_i)_{i\in I}$ de flèches de $\mathrm{Cat}$, \deleted{\unsure{si}} $f=\coprod f_i\in W$, les $f_i$ sont dans $W$\inserted{, $\forall i\in I$}.

% slanted
\margin{En fait, on \unsure{sait|soit} prouver l'implication en sens inverse, \deleted{\unsure{dans}} moyennant seulement Loc $(1,2,\deleted{\unsure{4}}5,6c')$ !}

Et si $i:Y_0\hookrightarrow Y$ est une flèche qui est à la fois imm. ouverte et imm. fermée (i.e. un iso. avec un \unsure{summand|sommand} direct de $Y$), alors $i$ est dans $\mathrm{Fib}_W$, \deleted{\unsure{ce qui à dire}}\inserted{\unsure{donc|soit} : si} \deleted{$f:X-Y$ est dans} $f:X_{\deleted{\unsure{0}}}\to\deleted{\unsure{X_2}}X'$ est une flèche dans $\mathrm{Cat}_{/Y}$, alors $f\in W$ implique $(f_0:X_0\to X'_0)\in W$, où $X_0=X\mid Y_0$, $X'_0=X'\mid Y_0$.

Mais il est assez scandaleux d'avoir à utiliser un axiome aussi fort, pour en conclure un résultat aussi anodin !

Je note ceci :

\uline{Proposition 9.4}. Soit $W\subset\mathrm{Fl}(\mathrm{Cat})$, satisfaisant Loc $(1)$ sans plus. Les conditions suivantes sont équivalentes (et entraînées par $L(1,2,3,\inserted{\unsure{4,}}5c')$

a) Si $f_1:X_1\to Y_1$, $f_2:X_2-Y_2$ sont deux flèches dans $W$, alors $f_1\amalg f_2\in W\iff f_1\in W,f_2\in W$

\deleted{\unsure{b) Pour que}}\unsure{b)} \deleted{\unsure{On a}} les deux \deleted{\unsure{propri}}conditions suivantes sur $W$ sont satisfaites :
\scanpage{249}
\mspage{123}
% archive sheet number: 248

$\underline{1}^{\circ}$) Si $f:X\unsure{-|\to}Y$ est dans $W$, alors il l'est couniversellement \unsure{sous|sur} $\varnothing$, i.e. pour \unsure{tt|$\forall$} $Z$ dans $\mathrm{Cat}$
\[
 f\amalg\mathrm{id}_Z:X\amalg Z\longrightarrow Y\amalg Z\quad\text{est dans }W.
\]

(\uline{NB} Ceci \deleted{\unsure{signifie que}}\inserted{équivaut à :} \unsure{tte somme finie|ttes sommes finies} de flèches $f_i\in W$ est dans $W$ --- et c'est $=$ pour sommes infinies, si Loc $(7)$, axiome des limites, est satisfait.)

$2^\circ$) Soit $Y_0\hookrightarrow Y$ une inclusion d'un \unsure{summand|sommand} direct d'un objet de $\mathrm{Cat}$. Alors si $X\xrightarrow{f}Y$ est dans $W$, de $=$ $X_0=X\mid Y_0-Y_0$.

(\uline{NB} Ceci équivaut à la condition que l'inclusion d'u\deleted{\unsure{ne}}\unsure{n} summand direct \deleted{dans $Y$} est toujours une $W$-fibration)

De plus, si $\pi_0(Y)$ est fini, ou si $W$ satisfait Loc $(7)$\inserted{(l'axiom. des limites)}, il suffit \deleted{dans $2^\circ$ \unsure{que}} d'exiger $2^\circ$) pour $Y_0$ une composante connexe de $Y$.

\inserted{c)} \deleted{\uline{Corollaire} 9.5. Soit $(f_i)_{i\in I}$ une} \inserted{Pour tte} famille \uline{finie} \inserted{$(f_i)_{i\in I}$} de flèches \deleted{de $\mathrm{Cat}$. Quand $W$ satisfait} $f_i:X_i-Y_i$ de $\mathrm{Cat}$, \deleted{Pour que} on a
\[
 \coprod f_i\in W\Longrightarrow f_i\in W\quad\text{pour tt }i\in I.
\]

De plus, si $W$ satisfait Loc $(7)$, alors il est inutile dans c) de se borner à $I$ fini.

(Et en tous cas, les conditions équivalentes a) b) c) impliquent pour \uline{tte} famille de flèches $(f_i)_{i\in I})$\unsure{;|,} \deleted{\unsure{que}} $\coprod f_i\in W\Rightarrow f_i\in W$ $\forall i\in I$.)
\scanpage{250}
\mspage{124}
% archive sheet number: 249

\deleted{Corollaire 9.5} \unsure{,}

D'après 9.3, ces conditions sont satisfaites moyennant Loc $(1,2,3,\inserted{4,}5c')$, mais $5c'$ est un marteau-pilon pour une mouche !

On aurait envie \deleted{\unsure{de f}} plutôt d'ajouter \deleted{\unsure{les propriétés}} un axiome aux axiomes Loc $(1)$ à Loc $(\deleted{\unsure{6}}7)$, qui avec Loc $(7)$ serait \unsure{l'axiome} le plus anodin de tous

% slanted
\margin{Il peut être commode \unsure{aussi} d'introduire les conditions encore plus faibles \uline{Loc} $(6\,\unsure{a})$ \deleted{Si $f_1,f_2$ \unsure{sont dans $W$}}, alors $f_1\amalg f_2$ \unsure{est} $\in W$. C'est en fait la condition \unsure{b $1^\circ$} de la prop.~9.4. \uline{Loc} $(6\,b)$ \unsure{Soit|Si} $f_1\amalg f_2$ est dans $W$, alors $f_1$ et $f_2$ \unsure{le sont}. $\Longleftrightarrow$ Soit $Y_0\subset Y$ un summand direct, alors $Y_0\hookrightarrow Y$ est \deleted{\unsure{asph}} dans $\mathrm{Fib}_W$.}

\uline{Loc} $(\deleted{\unsure{6|7}}5)$ \inserted{(Axiome de décomposition, forme faible)} : Si $(\unsure{f^n|f_n|f_i|f^i})_{\unsure{n|i}\in I}$ est une famille finie de flèches de $\mathrm{Cat}$, alors $f=\coprod f_i$ est dans $W$ * \unsure{ssi|ss} les $f_i$ le sont. (Ce sont donc les cond. équiv. de 9.4.)

\uline{Loc} $(\deleted{\unsure{6 bis|7 bis}}5\,\mathrm{bis})$ \inserted{(Axiome de décomp., forme forte)} : comme \uline{Loc} $(\unsure{8|5})$, mais sans supposer $I$ fini.

\uline{Loc} $(5)$ \quad a) Si $f_1,f_2\in W$, alors $f_1\amalg f_2\in W$
\par b) Si $f_1\amalg f_2\in W$ alors $f_1,f_2\in W$
\par ($\Longleftrightarrow$ b$'$) Soit $Y_0\hookrightarrow Y$ summ. direct,
\unsure{alors} $Y_0\hookrightarrow Y$ est dans $\mathrm{Fib}_W$)

\uline{Loc} $(5\,\mathrm{bis})$ comme Loc $(5)$ mais a) renforcé \inserted{(a bis)} pour une famille \inserted{(évent. infinie)} indexée par $I$ arbitraire \unsure{de flèches}.

Ainsi, moyennant Loc $(\deleted{\unsure{6}}7)$, les deux axiomes précédents sont équivalents. De plus, notons ceci :

\deleted{Proposition} \uline{Corollaire} 9.5. (à prop.~9.4.) L'axiome Loc \deleted{$(\unsure{6\,\mathrm{bis}})$}\inserted{$(5\,\mathrm{bis})$} équivaut \deleted{à la}\inserted{à chacune des} propriétés suivantes : \deleted{\unsure{Soit les}}

\textcircled{a} \inserted{Pour tte} $Y$ de $\mathrm{Cat}$, \unsure{décomposé|décomposée} en ses composantes connexes $Y_i$ ($i\in I=\pi_0(Y)$), \deleted{pour} et pour tte flèche $f:X\unsure{-|\to}Y$ dans $\mathrm{Cat}$, on a $f\in W$ \unsure{ss|ssi} les $f_i:X_i=X\mid Y_i-Y_i$ sont $\in W$.

\textcircled{b} Pour $Y,Y_i$ comme dans \textcircled{a}, pour tte
\scanpage{251}
\mspage{125}
% archive sheet number: 250

flèche $f:X\unsure{-|\to}X'$ dans $\mathrm{Cat}_{/Y}$, on a $f\in W$ ssi $\forall i\in I$, \deleted{\unsure{on a}} le morphisme $\unsure{f^i|f_i}:X_i-X'_i$ au dessus de $Y_i$ est dans $W$.

\uline{Corollaire 9.6}. Si l'axiome Loc $(\deleted{\unsure{6 bis}}5)$ est satisfait, alors l'axiome de connexité Loc $(8)$ équivaut à ceci : pour toute flèche $f:X-Y$ de $W$, avec $Y$ \unsure{$0$-connexe|$o$-connexe}, $X$ est $0$-connexe.

NB Les axiomes Loc $(8)$ (connexité) et Loc $(\deleted{\unsure{6}}5\,\mathrm{bis})$ \deleted{\unsure{ou Loc (8 bis)}} (de décomposition) ont l'air de nature voisine, \deleted{\unsure{je dirais $=$ les avantages,}} \deleted{\unsure{car}} étant liés l'un et l'autre \deleted{\unsure{ou que}} à la décomposition des objets de $\mathrm{Cat}$ en composantes connexes. (Néanmoins, je pense qu'ils sont indépendants.) Je présume qu'on n'aura jamais à travailler vraiment avec des $W$ qui ne satisfont à ces deux axiomes ! (En plus de Loc $(1,2,3)$, qui en constituent le ``minimum vital''.)

Mais apparemment il n'y en a pas besoin dans les développements théoriques généraux jusqu'ici. (Mais l'axiome de connexité Loc $(8)$ pourrait jouer un rôle, quand on parle de la \unsure{déterminabilité|détermination} réciproque des structures de \unsure{contractibilité|contractilité}, et de $W$-asphéricité, cf \unsure{P.S.})
\scanpage{252}
\mspage{126}
% archive sheet number: 251

Notons ceci :

% slanted
\margin{Cf version plus forte dans 9.\unsure{8|9}, plus bas (débarrassée de l'hyp. Loc (6)).}

\uline{Proposition 9.7}. Supposons \deleted{l'axiome} que $W$ satisfait $\mathrm{Loc}\,(1,2,3,\deleted{\unsure{6}}5)$, i.e. c'est un localiseur fondamental satisfaisant la condition de décomposition (forme faible). \deleted{\unsure{Le}} Si l'axiome Loc $(8)$ de connexité n'est \uline{pas} satisfait, alors \deleted{tous les \unsure{éléments} \illegible{1} \unsure{flèches}}\inserted{la sous-catégor. $\mathrm{Hot}^{\bullet}_W$} des $X\in\mathrm{Cat}$ non vides \unsure{est|\deleted{est}} équiv. à la cat. ponctuelle (donc $\mathrm{Hot}_W$ est équivalente soit, à la catégorie ponctuelle, \deleted{soit à cat} si

% slanted
\margin{\deleted{\unsure{Il suffit de savoir que la condition $2^\circ$ de b) soit satisfaite}}}

% slanted
\inserted{$\varnothing\to e$ est un iso dans $\mathrm{Hot}_W$, soit à $\Delta^1$ dans le cas contraire}

% slanted
\margin{Le premier cas est aussi celui où $\exists f\in W$ avec $\pi_0(f)$ \uline{non surjectif}, ou encore $\varnothing\to X$ dans $W$, avec $X$ non vide.}

Ceci \deleted{\unsure{généralise}}\inserted{est une variante du} th.~2 de \uline{XII} (p.~20), \deleted{\unsure{auquel on}} où, en faisant l'hypothèse Loc $(5c')$, plus forte que Loc $(\unsure{7|5})$ (et $=$\inserted{pour} Loc $(7\,\mathrm{bis})$), mais avec la conclusion plus forte que $W=\mathrm{Fl}(\mathrm{Cat})$

\uline{Dém} Soit
\[
 f:X\unsure{-|\to}Y,\quad f\in W,
\]
telle que
\[
 \pi_0(f)\quad\text{ne soit pas \uline{surjective}.}
\]

\deleted{Alors par Loc} Si $Y_0$ est une comp. connexe qui n'est pas dans $\mathrm{Im}\,\pi_0(f)$, alors $X_0=X\mid Y_0$ est vide, et \deleted{\unsure{par}}\unsure{par} Loc $(7)$ $\varnothing\to Y_0$ est dans $W$. Donc $\forall Z$ de $\mathrm{Cat}$, $\underset{\substack{\|\\\varnothing}}{\varnothing\times Z}\longrightarrow Y_0\times Z$ est dans $W$, \deleted{Pour toute}\inserted{et pour} \deleted{\unsure{toute}} flèche $g:Z-Z'$, $\mathrm{id}_{Y_0}\times g:Y_0\times Z\to Y_0\times Z'$ est dans $W$. Prenant $\unsure{g'|g}$ telle que $\pi_0(g)$ ne soit pas injective, p.ex. $e\amalg e\to e$, on \unsure{trouve|obtient} une flèche $\deleted{\unsure{g}}\unsure{g'|f}=\mathrm{id}_{Y_0}\times g$ \inserted{dans $W$}, telle que $\pi_0(\unsure{g|f})$ ne soit pas injective.
\scanpage{253}
\mspage{127}
% archive sheet number: 252

Soit donc $\deleted{\unsure{f|g}}\unsure{\varphi|f}:\deleted{X}I\to\deleted{Y}J$ dans $W$, avec $\pi_0(f)$ non injective. Procédant comme \unsure{tantôt|\deleted{tantôt}} en utilisant Loc $(\deleted{\unsure{6}}5)$, on est ramené au cas où de plus $\deleted{Y}J$ est $0$-connexe. (\uline{NB} dans les deux cas, on aurait pu aussi invoquer 9.6.) \deleted{On va utiliser les} Choisissons deux objets $\deleted{\unsure{x}}i_0,\deleted{\unsure{x}}i_1$ dans $\deleted{X}I$ dans deux composantes connexes distinctes, soient $\deleted{\unsure{y}}j_0,\deleted{\unsure{y}}j_1$ leurs images dans $\deleted{Y}J$. On va traiter $(I,i_0,i_1)$ et $(J,j_0,j_1)$ comme des ``intervalles homotopiques''.

Le fait que $i_0,i_1$ appartiennent à des composantes distinctes de $I$ implique que deux flèches $f_0,f_1:X\rightrightarrows Y$ sont ``homotopes par. à $(I,i_0,i_1)$'' i.e. il existe $h:X\times I\unsure{-|\to}Y$ dans $\mathrm{Cat}$ telle que $f_0=h\partial_0^X$, $f_1=h\partial_1^X$, où
\[
 \partial_0^X,\partial_1^X:X-I\times X
\]
sont donnés par $x\mapsto(\unsure{j|i}_\alpha\unsure{,|\times}x)$ $(\alpha\in\{0\,1\})$.

Pour trouver $h$, on décompose $\deleted{X}I$ en
\[
 \deleted{X}I=\deleted{X}I_0\amalg\deleted{X}I_1,\quad\text{avec } i_0\in I_0,\ i_1\in I,
\]
d'où
\[
 I\times \deleted{\unsure{I}}X\simeq(I_0\times \deleted{\unsure{I}}X)\amalg(I_1\times\unsure{X'|X)|X},
\]
et on définit $h$ par ses restrictions \inserted{$h_0,h_1$} aux deux morceaux $I_0\times X,I_1\unsure{<|\times}X$, savoir
\[
 h_0=\deleted{\unsure{\mathrm{pr}}}f_0\mathrm{pr}_2^{I_0},\quad
 h_1=f_1\circ\mathrm{pr}_2^{I_1},
\]
où
\[
 \mathrm{pr}_2^{I_\alpha}:I_\alpha\times X-X.
\]
\scanpage{254}
\mspage{1\deleted{\unsure{3}}2\deleted{\unsure{9|8}}8}
% archive sheet number: 253

Considérons maintenant le diagramme \inserted{(traits pleins)}

\diagram{\begin{tikzcd}[ampersand replacement=\&,row sep=large,column sep=small]
\deleted{X\amalg X}\arrow[d]\arrow[d,decorate,decoration={zigzag,segment length=3pt,amplitude=2pt},no head]
 \& X\vee X\arrow[dl,"{(\partial_0^X,\partial_1^X)}"']\arrow[dr,"{(\partial_0'{}^X,\partial_{\unsure{1|0}}'{}^X)}"] \& {} \\
I\times X\arrow[rr,"{\unsure{\varphi|\phi}\times\mathrm{id}_X}","{\substack{\|\\\deleted{\unsure{u}}\varphi_X}}"' ]
 \arrow[dr,"h"']\arrow[d,"\deleted{h}"']\arrow[d,decorate,decoration={zigzag,segment length=3pt,amplitude=3pt},no head]
 \& {} \& \deleted{\unsure{I}}J\times X\arrow[dl,dashed,"h'"] \\
{} \& Y \& {}
\end{tikzcd}}

, où $\partial'_0{}^X,\partial'_1{}^X$ sont définis comme $\partial_0^X,\partial_1^{\unsure{\alpha|X}}$, mais utilisant $j_0,j_1\in\mathrm{Ob}\,J$ au lieu de $i_0,i_1\in\mathrm{Ob}\,I$. S'il existait dans $\mathrm{Cat}$ une flèche $h'$ (pointillée) telle que $h'\circ\deleted{\unsure{u}}\varphi_X=h$, il s'ensuivrait que \deleted{les flèches} $h'$ est une homotopie de $f_0$ et $f_1$ suivant $(J,j_0,j_1)$.

Comme $J$ est \unsure{$0$-connexe|$o$-connexe}, il s'ensuivrait aussitôt que $f_0,f_1$ sont homotopes au sens habituel, dans $\mathrm{Cat}$, donc \deleted{\unsure{f}} que $f_0$ et $f_1$ ont $=$ image dans $\mathrm{Hot}_W$. Mais si on ne peut trouver une flèche $h'$ dans $\mathrm{Cat}$, du moins peut-on la trouver dans $\mathrm{Hot}_W$, puisque $\varphi_X\in W$ donc un iso dans $\mathrm{Hot}_W$ (et $h'$ dans $\mathrm{Hom}_W(J\times X,Y)$ \unsure{sera|en.} $=$ \unsure{uniquement|univoquement} déterminée). Donc dans $\mathrm{Hot}_W$ on a \deleted{\unsure{comme}}
\[
 f_0=h'\partial'_0{}^X,\qquad f_1=h'\partial'_1{}^X.
\]
\scanpage{255}
\mspage{129}
% archive sheet number: 254

Pour prouver que
\[
 f_0=f_1\quad\text{dans }\mathrm{Hom}_W(X,Y),
\]
il suffit de voir que \deleted{$\partial'_0{}^X$ et $\partial^{\unsure{''|'}}_1$}
\[
 \partial'_0{}^X=\partial'_1{}^X\quad\text{dans }\mathrm{Hom}_W(X,J\times X).
\]

Or en fait $\partial'_0{}^X$ et $\partial'_1{}^X$ sont homotopes, car
\diagram{\begin{tikzcd}[ampersand replacement=\&,column sep=small]
 e\arrow[r,shift left=.6ex,"j_0"]\arrow[r,shift right=.6ex,"j_1"'] \& J
\end{tikzcd}}
sont homotopes, \deleted{\unsure{d'un}} ($j_0$ et $j_1$ \unsure{appartiennent|appart.} à la $=$ comp. conn. de $J$), donc aussi \unsure{leur produit|leurs produits} par $\mathrm{id}_X$.

On voit donc que pour deux objets quelconques $X,Y$ de \deleted{$\mathrm{Cat}$, $\unsure{\mathrm{Hot}}$} $\mathrm{Cat}$, \unsure{toutes} les flèches de $X$ dans $Y$ sont égales dans $\mathrm{Hot}$. L'argument fait prouve $=$ que \deleted{toutes} les flèches \uline{dans $\mathrm{Hot}_W$} de $X$ dans $Y$ sont égales.

Cela signifie que \unsure{la} cat. $\mathrm{Hot}^{\unsure{*|\bullet}}_W$ (où \inserted{pour} deux objets quelc. \inserted{$X,Y$}, \deleted{\unsure{sont reliés par une flèche} \inserted{$\mathrm{Hom}(X,Y)$}} $\mathrm{Hom}_{\mathrm{Cat}}(X,Y)$ et \unsure{a|à} fortiori $\mathrm{Hom}_W(X,Y)$ est $\ne\varnothing$) est équivalente à la cat. ponctuelle.

Quant à $\mathrm{Hot}_W$ elle $=$, \deleted{qui est}\inserted{dont l'ens} des objets est $\{\varnothing\}\amalg\mathrm{Ob}\,\mathrm{Hot}^{*}_W$, \inserted{avec $\varnothing$ objet initial}, elle sera donc équivalente à la catégorie ponctuelle, ou bien à $\Delta^1$, selon que $\varnothing\to e$ est un iso dans $\mathrm{Hot}_W$, ou non.
\scanpage{256}
\mspage{130}
% archive sheet number: 255

% slanted
\margin{en fait, ces conditions sont $=$ en l'absence de l'hypothèse Loc $(\unsure{6|5})$, cf 9.\unsure{9|8} ci-dessous}

\uline{Commentaire} \unsure{9.8}. Si les axiomes Loc $(1,2,3,\deleted{\unsure{6}}5)$\unsure{$^{*}$} sont satisfaits, alors l'axiome \inserted{de connexité} Loc $(8)$ \uline{équivaut} à chacune des suivantes :

a) $\mathrm{Hot}_W$ n'est équivalente ni à $e$, ni à $\Delta^1$

b) $\mathrm{Hot}_W$ \deleted{n'est pas une cat. \unsure{formée}}\inserted{a une infinité de classes} d'isomorphie

c) \deleted{$\unsure{e\to}$} $f_0:X_0\to Y_0$ (flèche donnée d'avance \unsure{dans|\deleted{dans}} $\mathrm{Cat}$, telle que $\pi_0(f_0)$ n'est pas iso, \inserted{et $X_0\ne\varnothing$}) n'est \uline{pas} un iso. de $\mathrm{Hot}_W$.

d) $X_0$ et $Y_0$ ne sont pas isomorphes dans $\mathrm{Hot}_W$ (où \unsure{$X.$|$X_0$}, \unsure{$Y.$|$Y_0$} sont deux objets donnés d'avance de $\mathrm{Cat}$, avec $X_0,Y_0$ non vides et $\mathrm{card}\,\pi_0(X)\ne\mathrm{card}\,\pi_0(Y)$)

% double vertical rule beside the opening four lines
\margin{\rule[-3\baselineskip]{0.4pt}{4\baselineskip}\hspace{1.4pt}\rule[-3.15\baselineskip]{0.4pt}{4.15\baselineskip}}

On peut donc dire que l'axiome de connexité joue le rôle d'un \uline{axiome de non trivialité} pour la cat. $\mathrm{Hot}_W$ localisée de $\mathrm{Cat}$. Il sera satisfait automatiquement pour tous les localiseurs \inserted{fond.} $W$ qu'on sera amené à construire ou à introduire. \deleted{Donc}\inserted{Mais}, on n'aura pratiquement jamais à l'invoquer comme une hypothèse, en vue de certaines conclusions.

$\backslash$

Je m'aperçois finalement que dans la prop. 9.7 l'hypothèse Loc $(6)$ est inutile :

% slanted
\margin{\unsure{Cf mieux une variante|Cf une version} 9.12 (p.~13\unsure{\deleted{3}\inserted{4}|3} bis)}

% slanted
\margin{Loc $(6\unsure{\deleted{a}|a})$}

\uline{Prop. 9.\unsure{9|8}}. Soit $W$ un localiseur fondamental de $\mathrm{Cat}$ (i.e. on a Loc $(1,2,3)$ sans plus), \inserted{\uline{plus} la condition que la \uline{somme} de $f_1,f_2\in W$ est \deleted{dans} $\in W$, $^{*}$} où l'axiome de connexité n'est pas satisfait, alors $\mathrm{Hot}_W$ est équivalent à la catégorie ponctuelle (si $\varnothing\to e$ est un iso dans $\mathrm{Hot}_W$)$^{**}$ ou à $\Delta^1$ (dans le cas contraire), et $\mathrm{Hot}^{\bullet}_W$ ($\overset{\mathrm{df}}{=}\mathrm{Hot}_W\smallsetminus\{\varnothing\}$) est équivalente à la cat. ponctuelle.)

% slanted
\margin{* Cette condition est vérifiée notamment si $W$ satisfait de plus l'un des trois axiomes Loc $(4)$ (cf 9.\unsure{1|2}) ; Loc $(5c')$ (cf 9.3) ou Loc $(6\deleted{\unsure{a|bis}})$.}

$(**)$ Cela signifie aussi qu'il existe $Z\ne\varnothing$ dans $\mathrm{Cat}$, tel que $\varnothing\to Z$ soit iso dans $\mathrm{Hot}_W$.
\scanpage{257}
\mspage{131}
% archive sheet number: 256

La condition sur les sommes assure que dans $\mathrm{Hot}_W$, la somme de deux objets existe, et que $\mathrm{Cat}\to\mathrm{Hot}_W$ commute aux sommes de deux objets. Ceci dit, soit
\[ f:X\longrightarrow Y \]
une flèche telle que $f\in W$, mais $\pi_0(f)$ non surjectif. \inserted{Supposons d'abord $X\ne\varnothing$ (cas le moins évident).}

% slanted
\margin{\deleted{Si $X=\varnothing$, on prend \unsure{comme} dans} \unsure{cas}}

\deleted{\unsure{Si|donc}} $Y_0$ est un summand non vide de $Y$ tel que $f$ se factorise par $Y_1=Y\smallsetminus Y_0$, alors $f$ est \deleted{$f$ est} somme de
\[ f_{\unsure{1}}=X=X_1\longrightarrow Y_1\quad\text{et}\quad f_0:X_0=\varnothing\longrightarrow Y_0 \]
dans $\mathrm{Cat}$, donc aussi dans $\mathrm{Hot}_W$. Donc on a une somme $f=f_0\amalg f_1$ dans $\mathrm{Hot}_W$, qui est un iso., donc pour $\forall\,Z$ dans $\mathrm{Hot}_W$,

\deleted{\[ \mathrm{Hom}_W(Y_0,Z)\times\mathrm{Hom}_W(Y_1,Z) \]}
\[ \mathrm{Hom}_W(Y_0,Z)\times\mathrm{Hom}_W(Y_1,Z)
 \xrightarrow{\sim}\underbrace{\mathrm{Hom}_W(X_0,Z)}_{\substack{\parallel\\ e}}\times\mathrm{Hom}_W(X_1,Z) \]

(NB $\varnothing$ est objet initial de $\mathrm{Hot}_W$, l'étant de $\mathrm{Cat}$) i.e. le composé
\deleted{\[ Y_0(Z)\times Y_1(Z)\xrightarrow{\mathrm{pr}_2}Y_{\unsure{1}}(Z)\longrightarrow \]}
\[ Z(Y_0)\times Z(Y_1)\xrightarrow[\mathrm{pr}_2]{}Z(Y_1)\xrightarrow[Z(f_1)]{}Z(X_1) \]
est \unsure{bij.} Comme $X=X_1\ne\deleted{\unsure{\varnothing}}\varnothing$, donc aussi $Y_1\ne\varnothing$, et $Y_0\ne\varnothing$, \deleted{les \unsure{a}} les ens. $Z(Y_0),Z(Y_1),$ $Z(\unsure{\alpha_1|X_1})$ sont $\ne\varnothing$ si $Z\ne\varnothing$ ($\mathrm{Cat}$). \deleted{\unsure{Cela i}} Le composé \unsure{étant|\deleted{est} étant} bijectif, la première flèche $\mathrm{pr}_2$ est injective, ce qui implique
